% Solution: Milgrom & Roberts (1992) Ch. 8 Food for Thought 4 (public officials' pay). Input by \solitem.
% Check: check_T1.py (item 7).
\begin{exsol}{Milgrom \& Roberts Ch.~8, Food for Thought 4}{284--285}{Foundations}{2}
\asked Discuss the efficacy of proposals that tie pay increases for public officials to limits on their
outside earnings (lectures for pay; lobbying after leaving office).

\method The efficiency-wage model of Ch.~8 (pp.~250--255). An official does not sell favours iff
$p\,(w - \bar w)\,N \ge g$, where $w$ is pay in office, $\bar w$ the best alternative if dismissed, $p$ the
probability of detecting misconduct, $N$ the horizon multiplier and $g$ the one-time gain from
misconduct. The minimum deterrent pay is $w = \bar w + g/(Np)$ (formula 8.3). The exercise is discursive,
so this is an analytical answer.

\solution
\textbf{1. Where outside activities enter.} Outside income plays two roles.
(i) Lecture fees and later lobbying jobs from interest groups can be the \emph{payment for favours}: they
raise $g$, and they do not vanish if the official is found out, because the lobby values the same
connections, so the penalty $w - \bar w$ falls ($\bar w$ rises).
(ii) They are also a way of paying the official, which keeps the visible salary $w$ low and the rent
$w - \bar w$ of the job small. Both reduce $p(w-\bar w)N - g$.

\textbf{2. What the plan does.} Tying a pay rise to a cap on outside activities (a) raises $w$ and
(b) cuts $g$ and the outside component of $\bar w$. Both move the incentive constraint the right way.
In the model the plan works iff
\begin{itemize}
  \item the raise exceeds the outside income given up, otherwise $w - \bar w$ does not grow. With $p = 0.2$,
  $N = 4$, $\bar w = 100$ and $g = 80$ the deterrent salary is $w = 200$. An official who earns $30$ on the
  side needs only $170$, and a raise of $20$ in exchange for the ban does not deter;
  \item the limits are enforceable: $p$ must apply to the new rules (disclosure; bans on post-office
  lobbying with credible penalties), otherwise the ban is cosmetic;
  \item the post-office restriction is credible: a delayed lobbying job is a reward that cheating
  officials \emph{keep}, so only a binding cooling-off period reduces $g$.
\end{itemize}

\textbf{3. Costs and caveats.} (a) Participation: pay must still attract capable officials, whose
alternatives are lucrative. A cap lowers total pay unless the salary rises enough, which is politically
costly. (b) Not all outside activity is corrupt: lectures spread expertise, so a blunt ban destroys value.
(c) Payments can be displaced to relatives or consulting firms, so monitoring ($p$) is the weak link, as in
Macaulay's Lord Clive example (remote, powerful officials, low $p$). (d) A higher salary raises the rent
of honest officials too.

\textbf{4. Assessment.} The plan is effective as a package: a salary increase that at least replaces
the lost income, caps on activities that are payments for favours, and credible monitoring. A pay rise
without limits leaves $g$ untouched. Limits without a pay rise lower $w - \bar w$ and invite cheating or
exit.

\interp The proposal applies Clive's logic of liberal remuneration (p.~252) together with the monitoring
side of the model: pay and monitoring are substitutes in $p(w-\bar w)N \ge g$.

\verify \texttt{check\_T1.py} (item~7) computes the deterrent salary $200$, the salary $170$ with $30$ of
outside income, shows that a $+20$ raise against a $30$ ban fails to deter, and that a larger gain
($g = 120$) raises the required salary to $250$. The numbers are illustrative.

\answer{Deterrence needs $p(w - \bar w)N \ge g$. A package of a higher salary, limits on paid outside work
and enforceable post-office rules raises the rent and cuts the gain from favours, so it can work. It fails if
the raise is smaller than the income given up, if the limits cannot be monitored, or if the cooling-off
rule is not credible.}
\end{exsol}
